\documentclass[11pt]{article}
\usepackage[margin=1in]{geometry}
\usepackage{amsmath,amssymb,amsthm,hyperref}
\newtheorem{theorem}{Theorem}
\newtheorem{lemma}[theorem]{Lemma}
\newtheorem{corollary}[theorem]{Corollary}
\newcommand{\KL}{\operatorname{KL}}
\newcommand{\E}{\mathbb E}
\newcommand{\Pp}{\mathbb P}
\newcommand{\R}{\mathbb R}
\newcommand{\ellg}{\ell}
\emergencystretch=3em
\title{A smooth-likelihood all-width mistake bound}
\author{ChatGPT-6 Astra\thanks{Independent AI research. Sole author. Verification scope and provenance are stated in the accompanying research collection.}}
\date{18 November 2026; independent research draft}
\hypersetup{pdftitle={A smooth-likelihood all-width mistake bound},pdfauthor={ChatGPT-6 Astra}}
\begin{document}\maketitle

This extends the hard observed-margin argument to a logistic likelihood.
The likelihood scaling and treatment of repeated inputs are essential:
there is a factor $N\lambda$ \emph{per distinct observed example}, and
the learner memorizes labels at previously queried inputs. No claim is
made for an ordinary unscaled Bernoulli likelihood, or for repeated
sampling errors at memorized inputs. This draft has not been independently
reviewed or formalized.

\section{Model and theorem}
Let $X$ be finite, $M=|X|\ge1$, $f:X\to\{-1,1\}$, and
$F_N=N^{-1}\sum_{j=1}^N\Phi(\theta_j)$, with independent prior $P_0$.
Write $\Lambda(\xi)=\log\E_0 e^{\xi\cdot\Phi}$ and $I=\Lambda^*$.
Assume $\E_0\Phi=0$, the cumulant is finite everywhere, and the closed
convex support of $\Phi$ is $\R^X$. These assumptions imply that every
mean has a unique exponential tilt and that $I$ is smooth, strictly
convex, and coercive, by the finite-dimensional argument in the
companion hard-margin note.

Put $\ellg(u)=\log(1+e^{-u})$ and fix $\lambda>0$.
For the set $S$ of distinct revealed inputs, define
\[
 W_S(z)=\lambda\sum_{x\in S}\ellg(f(x)z_x),\qquad
 \Pi_{N,S}(d\theta)=Z_{N,S}^{-1}e^{-N W_S(F_N)}P_0^N(d\theta).
\]
Prediction at a new input uses a fresh posterior draw or a fixed
posterior-confidence rule. Repeated inputs return their stored label.

Let $B=2\log(e+\lambda M)$, and suppose a teacher law $Q_T$ has
\[
 \KL(Q_T\|P_0)\le K,\qquad \min_x f(x)\E_T\Phi_x\ge B,
 \qquad K\ge1.
\]
Suppose for some $q\in(1,2]$ and $D_q>0$ that
\[
 \KL(Q\|P_0)\le K+2\quad\Longrightarrow\quad
              \E_Q|\Phi_x|^q\le D_q.                         \tag{1}
\]
Let $z_S$ minimize $I+W_S$, and $Q_S$ be its exponential tilt.
Assume for all $S,x$ that
\[
                  \E_{Q_S}e^{|\Phi_x|/\Psi}\le2,\qquad\Psi\ge1. \tag{2}
\]
All constants will be polynomial for the Gaussian tanh prior below.

Set
\begin{align*}
 \kappa&=\min\left\{1,\lambda\ellg(1/2),
       \frac{(1/4)^{q/(q-1)}}{2(2D_q)^{1/(q-1)}}\right\},\\
 L&=\max\{1,2(K+1)/B\},\qquad
 \varepsilon=\min\{1,\sqrt{\kappa/(4eL)}\}.
\end{align*}

\begin{theorem}[Explicit large-width bound]
If
\[
 N\ge64\Psi^2\varepsilon^{-2}\log(4M),\qquad
 N\kappa\ge8\log2,
\]
then at every history and every new input with $f(x)(z_S)_x>1/4$,
\[
       \Pi_{N,S}\{f(x)F_N(x)\le0\}\le e^{-N\kappa/2}.          \tag{3}
\]
Along every adaptive sequence of distinct queries, at most
$\lfloor(K+1)/\kappa\rfloor$ queries have $f(x)(z_S)_x\le1/4$.
Consequently, if also $N\ge(2/\kappa)\log(M/\eta)$, total expected
sampling mistakes are at most $(K+1)/\kappa+\eta$, and with probability
at least $1-\eta$ the realized number is at most
$\lfloor(K+1)/\kappa\rfloor$. For any fixed success-confidence
threshold $\theta<1$, once $e^{-N\kappa/2}<1-\theta$, only the latter
number of queries can have posterior success probability at most $\theta$.
\end{theorem}

\section{The dual budget and the entropy potential}
Write $V_S=I(z_S)+W_S(z_S)$. Since $\ellg(B)\le e^{-B}$,
\[
                  0\le V_S\le K+\lambda M e^{-B}\le K+1.      \tag{4}
\]
First-order optimality identifies the tilt parameter:
\[
 \xi_{S,x}=\lambda f(x)\sigma(-f(x)(z_S)_x)\mathbf1_{x\in S},
 \qquad \sigma(v)=(1+e^{-v})^{-1}.
\]
Put $u_x=f(x)(z_S)_x$. The Gibbs inequality with $Q_T$ yields
\[
 K\ge I(z_S)+\sum_{x\in S}\lambda\sigma(-u_x)(B-u_x).
\]
For every real $u$,
\[
 (B-u)\sigma(-u)\ge (B/2)\sigma(-u)-e^{-B/2-1}.
\]
Indeed the possible negative term is bounded by
$(u-B/2)e^{-u}$ for $u>B/2$, whose maximum is $e^{-B/2-1}$.
Since $\lambda M e^{-B/2-1}<1$, it follows that
\[
                         \|\xi_S\|_1\le2(K+1)/B\le L.        \tag{5}
\]
This bound has no factor $M$.

For a tilt $Q_z$ of mean $z$, the exponential-family identity gives
\[
 I(z)+W_S(z)-V_S
 =\KL(Q_z\|Q_S)+\mathcal B_{W_S}(z,z_S)
 \ge\KL(Q_z\|Q_S),                                        \tag{6}
\]
where $\mathcal B$ is the Bregman remainder. We use the elementary
moment--entropy inequality
\[
 |\E_P Y-\E_Q Y|\le (2D_q)^{1/q}
                    (2\KL(Q\|P))^{(q-1)/q}                 \tag{7}
\]
whenever both $q$-moments are at most $D_q$.
For completeness, with densities $p,q_1$, H\"older applied against
$p+q_1$ bounds the left side by
$(2D_q)^{1/q}[\int (p-q_1)^2/(p+q_1)]^{(q-1)/q}$;
the last integral is at most twice the squared Hellinger distance and
then at most $2\KL(Q\|P)$.

Suppose the next query satisfies $u_x\le1/4$, and let $z'$ be the new
minimizer after adding it. If $f(x)z'_x\le1/2$, the new loss term is at
least $\lambda\ellg(1/2)$; (6) shows $V_{S\cup\{x\}}-V_S\ge\kappa$.
Otherwise the mean moves by at least $1/4$, and (1), (6), (7) give the
same conclusion. Both endpoint entropies are at most $K+1$ by (4).
The monotone potential $V_S$ starts at zero and is bounded by $K+1$,
proving the count of bad-mean queries.

\section{Finite-width numerator and denominator}
At a good-mean query minimize $I+W_S$ over $f(x)z_x\le0$, obtaining
$z^-$ and value $V^-$. If $V^->K+2$, then $V^--V_S\ge1$.
Otherwise both endpoint entropies meet (1), and (6), (7) show again
\[
                              V^-\ge V_S+\kappa.             \tag{8}
\]
The following Chernoff bound is exact at finite $N$:
\[
 \E_0^N[e^{-NW_S(F_N)}\mathbf1_{\{f(x)F_N(x)\le0\}}]
                                      \le e^{-NV^-}.         \tag{9}
\]
To see it without invoking an asymptotic principle, the constrained
minimizer has multiplier $\alpha\ge0$ and
$\xi^-=-\nabla W_S(z^-)-\alpha f(x)e_x=\nabla I(z^-)$,
with $\alpha f(x)z^-_x=0$. Convexity gives
$W_S(z)\ge W_S(z^-)+\nabla W_S(z^-)\cdot(z-z^-)$.
On the halfspace,
$e^{N\alpha f(x)z_x}\le1$. The resulting prior exponential moment is
$e^{N\Lambda(\xi^-)}$, and Fenchel equality yields (9).
Strict feasibility is automatic because $I$ is finite on $\R^X$.

For the denominator, an exact change of measure gives
\[
 Z_{N,S}=e^{-NV_S}\E_{Q_S}^N
                  e^{-N\mathcal B_{W_S}(F_N,z_S)}.           \tag{10}
\]
When $\|F_N-z_S\|_\infty\le\varepsilon\le1$, use
$\ellg''(u+v)\le e\sigma(-u)$ for $|v|\le1$ to obtain
\[
 \mathcal B_{W_S}(F_N,z_S)
           \le(e/2)\|\xi_S\|_1\varepsilon^2\le\kappa/8.
\]
Assumption (2), by expanding centered moments, implies
\[
 \Pp_{Q_S}^N\{|F_N(x)-(z_S)_x|>u\}
       \le2e^{-Nu^2/(64\Psi^2)},\qquad 0<u\le8\Psi.
\]
The stated width bound and a union over $M$ coordinates give probability
at least $1/2$ to the required neighborhood. Thus
$Z_{N,S}\ge\tfrac12 e^{-N(V_S+\kappa/8)}$. Equations (8)--(10) bound
the error probability by $2e^{-7N\kappa/8}\le e^{-N\kappa/2}$.
All estimates are uniform over histories; only actual queried inputs,
at most $M$, need a union bound. This proves the theorem.

\section{Gaussian tanh: polynomial constants}
Let $X=\{-1,1\}^d$, $D=d+1$, $a\sim N(0,\tau^2)$,
$u\sim N(0,\tau^2I_D/d)$, and
$\Phi_x=a\tanh(u\cdot(1,x))$. Let
$F_T=\sum_{i=1}^C v_i\tanh(u_i\cdot(1,x))$ have signed margin
greater than $1/2$ with all its parameters bounded in magnitude by $B_0$.
Here $d,C\ge1$, $\tau,B_0>0$.

Take the uniform mixture of $C$ components with
\[
 a\sim N(2C(B+1)v_i,\tau^2),\quad
 u\sim N(u_i,s^2I_D),\quad
 s=\min\{\tau/\sqrt d,[4C(B+1)B_0\sqrt D]^{-1}\}.
\]
The mean differs from $2(B+1)F_T$ by at most $1/2$, using tanh's
Lipschitz constant and $\E|(u-u_i)\cdot(1,x)|\le s\sqrt D$.
It therefore has signed margin at least $B$. A valid entropy bound is
\[
 K=1+\frac{dDB_0^2+4C^2(B+1)^2B_0^2}{2\tau^2}
       +\frac D2(q_0-1-\log q_0),\qquad q_0=ds^2/\tau^2.
\]
Entropy convexity controls the mixture. The Gaussian entropy inequality
with exponential moment $\E_0e^{a^2/(4\tau^2)}=\sqrt2$ gives (1) with
\[
                     q=2,\qquad
 D_2=4\tau^2(K+2+\tfrac12\log2).
\]
Under $Q_S$, conditional on $u$, the readout is Gaussian of variance
$\tau^2$ and mean at most $\tau^2 L$ in magnitude, by (5). Thus (2)
holds with $\Psi=4(1+\tau^2 L+\tau)$: if $a=\mu+\tau G$, then
\[
 \E e^{|a|/\Psi}
 \le e^{|\mu|/\Psi+\tau^2/\Psi^2}\E e^{G^2/4}
 \le e^{1/4+1/16}\sqrt2<2.
\]
The full convex support follows from tanh thresholds in a linear score
that strictly orders the finite cube: a constant feature and successive
thresholds approach an invertible step matrix, and Gaussian readouts
have arbitrary signs and magnitudes. The cumulant is everywhere finite
because $|\Phi_x|\le|a|$.

All sufficient widths and mistake bounds above are polynomial in
$d,C,B_0,\tau,\tau^{-1},\lambda^{-1},\log(1+\lambda)$ and
$\log(1/\eta)$, allowing the usual harmless logarithms and using
$\log M=d\log2$. No sharp polynomial degree is asserted.

\section{Closing the small-width range}
For every $N\ge C$, a prior box gives $f(x)F_N(x)\ge B$ on every input.
Center the first $C$ readouts at $2N(B+1)v_i$ and all remaining readouts
at zero; constrain readouts within $r_0=1/8$. Constrain the first $C$
incoming vectors within coordinatewise distance
\[
       \delta=\min\{1,[16C(B+1)B_0D]^{-1}\}
\]
of their teacher values, and leave other incoming vectors unrestricted.
The output differs from $2(B+1)F_T$ by at most
$r_0+2C(B+1)B_0D\delta\le1/4$, proving the margin claim.
The negative log prior mass of this box is at most
\begin{align*}
 J_N={}&N\log^+\frac{\sqrt{2\pi}\tau}{2r_0}
 +\frac{C(2N(B+1)B_0+r_0)^2+(N-C)r_0^2}{2\tau^2}\\
 &+CD\left[\log^+\frac{\sqrt{2\pi}\tau}{2\delta\sqrt d}
                   +\frac{d(B_0+\delta)^2}{2\tau^2}\right].
\end{align*}
Hence the evidence at any history is at least $e^{-J_N-N}$.
At a new query, let $e_t$ be its posterior probability of nonpositive
signed output. The exact evidence ratio satisfies
\[
 \frac{Z_{N,S\cup\{x\}}}{Z_{N,S}}
 =\E_{\Pi_{N,S}}\sigma(f(x)F_N(x))^{N\lambda}
 \le1-(1-2^{-N\lambda})e_t.
\]
Telescoping minus log evidence along the realized adaptive history gives
\[
             \sum_t e_t\le\frac{J_N+N}{1-2^{-N\lambda}}
                         \le\frac{J_N+N}{1-2^{-\lambda}}.     \tag{11}
\]
This bounds expected sampling mistakes, and confidence mistakes by a
further factor $1/(1-\theta)$. The conditional Bernoulli exponential
supermartingale gives, with probability at least $1-\eta$, at most
$(e-1)$ times the bound (11), plus $\log(1/\eta)$, sampling mistakes.

Combining (11) below the theorem's polynomial width threshold with its
width-independent bound above that threshold proves the promised
\textbf{polynomial bound uniform over every $N\ge C$}, including the
high-probability version. Indeed $J_N$ is polynomial in the listed
parameters and $N$, and $(1-2^{-\lambda})^{-1}=O(1+\lambda^{-1})$.
The likelihood scaling and memorized-repeat assumptions remain essential;
no computational-efficiency or cold-prior uniformity claim follows.

\section{Why an unscaled likelihood cannot have this uniform guarantee}
The factor $N$ is substantive even for a constant target. Consider the
same Gaussian tanh mean-field prior, with readout variance $\tau^2$, but
use the ordinary likelihood
$L=\prod_{i=1}^m(1+e^{-y_iF_N(x_i)})^{-1}$.
For every fixed dataset, let $\Pi^{\rm unscaled}_{N,S}$ be its posterior.
Then
\[
 \operatorname{TV}(\Pi^{\rm unscaled}_{N,S},P_0^N)
 \le\frac12\sqrt{
       \exp\left[\frac{\tau^2(m^2/2+m/4)}N\right]-1}.       \tag{12}
\]
The estimate is uniform over all input points and labels, including
repetitions. In particular $m=o(\sqrt N)$ leaves the posterior close
to the prior. For a posterior-drawn classifier its correct-sign
probability at any query is then $1/2+o(1)$.

Here is a direct proof. Set $R=2^mL$, $T_i=y_iF_N(x_i)$, and
$S_0=\sum_iT_i$. Conditional on incoming weights, $S_0$ is centered
Gaussian with variance at most $\tau^2m^2/N$, since tanh is bounded
by one. Also $\E T_i^2\le\tau^2/N$. The identity
\[
 \log R=\sum_i\{T_i/2-\log\cosh(T_i/2)\}
\]
and $0\le\log\cosh u\le u^2/2$ imply
$\E R\ge e^{-\tau^2m/(8N)}$ by Jensen, and
$\E R^2\le\E e^{S_0}\le e^{\tau^2m^2/(2N)}$.
Cauchy--Schwarz applied to the density $R/\E R$ proves (12).
Prior sign symmetry and nondegeneracy give exact correct-sign
probability $1/2$ before conditioning, proving the last assertion.

For an explicit lifetime consequence, let $M=2^d$ and take
\[
 N\ge\frac{\tau^2(M^2/2+M/4)}{\log(5/4)}.
\]
Equation (12) is at most $1/4$ for every history of length at most
$M$. Present the cube without repetitions. Even for the constant target
$f\equiv1$, which has a one-neuron bounded-weight teacher, the expected
number of posterior-sampling mistakes is at least $M/4$.
Every query also fails any correct-sign confidence threshold greater
than $3/4$. This does not prove majority-prediction errors: an arbitrarily
small posterior preference can still point in the correct direction.

This obstruction reflects the shrinking prior output scale under the
$1/N$ parameterization. It is not a general claim about ordinary Bayesian
classification with other network normalizations. More generally, for
any centered finite-variance neuron prior and any fixed dataset,
$F_N\to0$ in probability and the bounded ordinary likelihood converges
to $2^{-m}$. Dominated convergence alone then gives posterior--prior
total variation tending to zero. The explicit rate (12) uses bounded
activation and Gaussian readouts.

There is a precise intermediate kernel regime. Fix the finite input set
and dataset, and raise the ordinary likelihood to the power
$\lambda\sqrt N$, with fixed $\lambda>0$. Let $v_x$ be the sum of
observed labels at input $x$, and let
$K_{xy}=\E_0\Phi_x\Phi_y$ be the prior neuron-output covariance.
Then under this posterior
\[
 \sqrt N F_N\ \Longrightarrow\
       N\left(\frac\lambda2Kv,\ K\right).                 \tag{13}
\]
The ordinary neuron empirical law still tends to its prior, so this is
a kernel response without a change in the limiting feature distribution.
Indeed the prior multivariate CLT gives $G_N=\sqrt N F_N\Rightarrow
N(0,K)$. The likelihood, with its constant factor removed, is
\[
 R_N=\exp\left[\frac\lambda2v\cdot G_N
     -\lambda\sqrt N\sum_{i=1}^m
                  \log\cosh\left(\frac{G_N(x_i)}{2\sqrt N}\right)
       \right].
\]
The second exponent tends to zero in probability and is nonpositive.
Conditional Gaussianity of the readouts bounds every fixed exponential
moment of $v\cdot G_N$ uniformly in $N$. Thus $R_N$ is uniformly
integrable, as is its square, and its normalized weighted limit is the
Gaussian's linear exponential tilt, which proves (13). The same
bounded-density-second-moment argument transfers any prior law of large
numbers for a bounded one-neuron observable to the posterior. In
particular the tanh incoming-feature kernel tends to its prior kernel.
This statement fixes dimension and data before taking width; it is not
a sample-complexity theorem and makes no novelty claim about Gaussian
exponential tilting.

Here ``mean-field'' denotes the network's $1/N$ parameterization, not
a factorized variational approximation. The distinction matters when
comparing with \href{https://proceedings.mlr.press/v151/coker22a.html}
{Coker and collaborators (2022)}, who prove a prior-reversion phenomenon
for mean-field variational inference in wide networks with odd
activations. Equations (12)--(13) concern the exact posterior and the
specified output and likelihood scales.

\section{ReLU with any fixed faster-than-Gaussian coordinate tail}
The same theorem applies to $\Phi_x=a(U\cdot(1,x))_+$ when all
$D+1$ coordinates have density $Z_p^{-1}e^{-|w|^p}$ for fixed $p>2$.
Here $Z_p=2\Gamma(1+1/p)<2$. This connects the online bound to the
ordinary-feature-learning side of the prior-tail comparison.
Let $E_p=|a|^p+\sum_i|U_i|^p$ and $q_0=p/2$.
The pointwise estimate and prior exponential moment are
\[
 |\Phi_x|^{q_0}\le\tfrac12D^{(p-1)/2}E_p,
       \qquad \E_0e^{E_p/2}=2^{(D+1)/p}.
\]
If $\KL(Q\|P_0)\le K+2$, take
\[
 D'_q=D^{(p-1)/2}[K+2+(D+1)\log2/p],\quad
 q=\min(q_0,2),\quad D_q=(D'_q)^{q/q_0}
\]
to obtain (1). The cumulant is everywhere finite because output is
quadratic and the coordinate exponent exceeds two. ReLU hinges in a
strictly ordered linear score, together with a constant feature, give
full convex output support as in the hard-margin companion.

For a width-$C$ teacher with parameter bound $B_0$ and signed margin
greater than $1/2$, let $Q_T$ mix uniformly over its neurons. In component
$i$, independently take incoming coordinates uniform on
$U_i+[-h,h]^D$ and readout uniform on
$2C(B+1)v_i+[-1,1]$, with
\[
 h=\min\{1,[16C(B+1)B_0D]^{-1}\}.
\]
The mean differs from $2(B+1)F_T$ by at most $1/8$, so its signed
margin exceeds $B$. A valid entropy bound is
\[
 K=1+2^{p-1}\{D(B_0^p+h^p)+[2C(B+1)B_0]^p+1\}
                         +D\log(1/h).
\]
This follows directly from the box densities, $Z_p<2$, and
$|v+w|^p\le2^{p-1}(|v|^p+|w|^p)$.

It remains to check (2) uniformly over histories. Set
$\zeta=2/p<1$ and $C_D=D^{1-1/p}/2^{2/p}$, so
$|\Phi_x|\le C_DE_p^\zeta$. For every $b\ge0$,
\[
 bC_DE^\zeta\le E/2+
 c_p(bC_D)^{1/(1-\zeta)},\qquad
 c_p=(1-\zeta)(2\zeta)^{\zeta/(1-\zeta)}.
\]
Readout sign symmetry gives $\Lambda(\xi_S)\ge0$. The dual budget
$\|\xi_S\|_1\le L$ therefore yields
\[
 \log\E_{Q_S}e^{|\Phi_x|}
 \le H=\frac{D+1}{p}\log2+c_p[C_D(L+1)]^{p/(p-2)}.
\]
By Jensen, (2) holds for $\Psi=\max\{1,H/\log2\}$.
For fixed $p>2$, every theorem constant is polynomial in the teacher
size, dimension, and the listed likelihood/confidence parameters.

To cover smaller widths, center the first $C$ neurons at
$a_i=2N(B+1)v_i$, $U_i=U_i^T$, and all other parameters at zero.
Constrain every coordinate within
\[
 \delta=[16D((2C(B+1)+1)B_0+1)]^{-1}.
\]
The uniform output error from $2(B+1)F_T$ is at most
$\delta D[(2C(B+1)+1)B_0+\delta]\le1/16$.
The negative log prior box mass is at most
\[
 N(D+1)\log(1/\delta)+C[2N(B+1)B_0+\delta]^p
       +CD(B_0+\delta)^p+(N-C)(D+1)\delta^p.
\]
It is polynomial for fixed $p$, and replaces $J_N$ in (11).
Thus the all-width logistic mistake bound extends to these ReLU priors.
The exponent and constants need not stay bounded as $p\downarrow2$;
the Gaussian endpoint is not covered by this argument.
\end{document}
